% Ad Atticum 3.22 (ad-atticum-03-22)
% Source: https://raw.githubusercontent.com/PerseusDL/canonical-latinLit/master/data/phi0474/phi057/phi0474.phi057.perseus-lat1.xml
% Fetched by scripts/fetch_latin.py

% Dateline (Perseus): Scr. partim Thessalonicae partim Dyrrachi vi K. Dec. a. 696 (58).

\ciceroLetterOpener{CICERO ATTICO salutem}

\ciceroSection{1} etsi diligenter ad me Quintus frater et Piso quae essent acta scripserant, tamen vellem tua te occupatio non impedisset quo minus, ut consuesti, ad me quid ageretur et quid intellegeres perscriberes. me adhuc Plancius liberalitate sua retinet iam aliquotiens conatum ire in Epirum. spes homini est iniecta non eadem quae mihi posse nos una decedere; quam rem sibi magno honori sperat fore. sed iam, cum adventare milites dicentur, faciendum nobis erit ut ab eo discedamus. quod cum faciemus, ad te statim mittemus, ut scias ubi simus.

\ciceroSection{2} Lentulus suo in nos officio, quod et re et promissis et litteris declarat, spem nobis non nullam adfert Pompei voluntatis; saepe enim tu ad me scripsisti eum totum esse in illius potestate. de Metello scripsit ad me frater quantum speraret profectum esse per te.

\ciceroSection{3} mi Pomponi, pugna ut tecum et cum meis mihi liceat vivere et scribe ad me omnia. premor luctu, desiderio cum omnium rerum tum meorum qui mihi me cariores semper fuerunt. cura ut valeas.

\ciceroSection{4} ego quod per Thessaliam si irem in Epirum perdiu nihil eram auditurus et quod mei studiosos habeo Dyrrachinos, ad eos perrexi, cum illa superiora Thessalonicae scripsissem. inde cum ad te me convertam, faciam ut scias, tuque ad me velim omnia quam diligentissime cuicuimodi sunt scribas. ego iam aut rem aut ne spem quidem exspecto. data vi Kal. Decembr. Dyrrachi.
